%% ============================================================
%%  Capítulo 5 — Riscos e Mitigações
%% ============================================================
\chapter{Riscos e Mitigações}
\label{chap:riscos}

\section{Identificação de Riscos}

A análise de segurança operacional identificou vários riscos e implementou mitigações correspondentes.

\subsection{Matriz de Riscos}

\begin{table}[h!]
\centering
\caption{Riscos identificados e mitigações implementadas}
\label{tab:riscos}
\rowcolors{2}{gray!15}{white}
\begin{tabularx}{\textwidth}{|X|X|}
\hline
\textbf{Risco} & \textbf{Mitigação} \\ \hline
Token inválido ou forjado & Validação RS256 via \gls{jwks}, verificação de \texttt{issuer} e \texttt{audience} \\ \hline
Acesso indevido a recursos & \gls{rbac} no \textit{backend} com \texttt{require\_role()} por endpoint \\ \hline
Administrador sem segundo fator & \gls{totp} obrigatório e validação de \texttt{amr} na \gls{api} \\ \hline
Sessões ativas após mudança de papel & Revogação explícita (\texttt{DELETE /users/\{id\}/sessions}) no Mover e Leaver \\ \hline
Credenciais expostas no código & Variáveis de ambiente via \texttt{.env} (não versionado) \\ \hline
Falta de rastreabilidade & Eventos Keycloak (PostgreSQL) e endpoint \texttt{/admin/audit} protegido \\ \hline
\end{tabularx}
\end{table}

\section{Mitigações Implementadas}

\subsection{1. Token Inválido ou Forjado}

\textbf{Risco:} Um atacante consegue criar um token falso ou modificar um token legítimo e usá-lo para aceder a recursos protegidos.

\textbf{Mitigação:}
\begin{itemize}
  \item Todos os tokens são validados com a chave pública RS256 do Keycloak
  \item Verifica-se a assinatura HMAC-SHA256 contra a chave pública
  \item Verifica-se que o \texttt{issuer} corresponde ao Keycloak esperado
  \item Verifica-se que o \texttt{audience} é o cliente FastAPI
  \item Verifica-se que o token não expirou
\end{itemize}

\subsection{2. Acesso Indevido a Recursos}

\textbf{Risco:} Um utilizador consegue aceder a endpoints que não deveria, contornando o controlo de acesso.

\textbf{Mitigação:}
\begin{itemize}
  \item Cada endpoint está protegido por um decorador \texttt{require\_role()} que valida o papel do utilizador
  \item Visitantes não podem aceder a \texttt{/colaborador/*} ou \texttt{/admin/*}
  \item Colaboradores não podem aceder a \texttt{/admin/*}
  \item A validação ocorre na FastAPI, independentemente do Keycloak
\end{itemize}

\subsection{3. Administrador sem Segundo Fator}

\textbf{Risco:} Um utilizador com papel de administrador consegue aceder a áreas críticas sem passar por validação de segundo fator, mesmo que tenha credenciais válidas.

\textbf{Mitigação:}
\begin{itemize}
  \item A \textit{Required Action} \texttt{CONFIGURE\_TOTP} é ativa para administradores
  \item No primeiro login, o Keycloak força a configuração de uma aplicação TOTP
  \item Cada login de administrador requer inserir um código de 6 dígitos
  \item O token inclui a \textit{claim} \texttt{amr: ["otp"]} apenas se \gls{mfa} bem-sucedido
  \item A \gls{api} valida \texttt{amr} antes de permitir acesso a \texttt{/admin/mfa-area}
\end{itemize}

\subsection{4. Sessões Ativas após Mudança de Papel}

\textbf{Risco:} Após um utilizador ser movido (papel alterado) ou desativado, as suas sessões antigas permanecem válidas durante o tempo de vida do token (300 segundos), permitindo acesso indevido.

\textbf{Mitigação:}
\begin{itemize}
  \item No fluxo Mover, a \gls{api} revoga todas as sessões do utilizador após alterar papel/grupo
  \item No fluxo Leaver, a \gls{api} revoga todas as sessões antes de desativar a conta
  \item O utilizador é forçado a fazer reautenticação para obter um novo token
  \item Não há janela de acesso indevido
\end{itemize}

\subsection{5. Credenciais Expostas no Código}

\textbf{Risco:} Senhas, segredos de cliente e chaves privadas ficam expostos no repositório git ou são incluídas no código fonte.

\textbf{Mitigação:}
\begin{itemize}
  \item Todas as credenciais são carregadas via variáveis de ambiente
  \item O ficheiro \texttt{.env} contém os segredos e \textbf{não} é versionado (incluído em \texttt{.gitignore})
  \item O repositório contém apenas \texttt{.env.example} como template
  \item Não há \textit{hardcoding} de passwords no código Python
  \item O \texttt{realm-export.json} contém hashes de passwords, não senhas em texto limpo
\end{itemize}

\subsection{6. Falta de Rastreabilidade}

\textbf{Risco:} Não há registo de quem fez o quê, quando, impossibilitando a detecção de atividades suspeitas ou a investigação de incidentes.

\textbf{Mitigação:}
\begin{itemize}
  \item Todos os logins (sucesso e falha) são registados no Keycloak
  \item Todas as operações de administração (\gls{jml}) são registadas
  \item Os eventos persistem na base de dados PostgreSQL
  \item O endpoint \texttt{/admin/audit} permite consultar eventos com filtragem por tipo
  \item O dashboard inclui uma aba de Auditoria com visualização e estatísticas
\end{itemize}

\section{Testes de Segurança}

\subsection{Testes Unitários}

Quatro testes automatizados validam a lógica de segurança:

\begin{table}[h!]
\centering
\caption{Testes unitários de segurança}
\label{tab:testes}
\begin{tabularx}{\textwidth}{|l|X|}
\hline
\textbf{Teste} & \textbf{O que valida} \\ \hline
\texttt{test\_require\_role\_allows\_admin} & \texttt{require\_role("admin")} aceita token com papel \texttt{admin} \\ \hline
\texttt{test\_require\_role\_denies\_non\_admin} & Papel incorreto retorna HTTP~403 \\ \hline
\texttt{test\_require\_mfa\_denies\_without\_otp} & Ausência de \texttt{amr=otp} retorna HTTP~403 \\ \hline
\texttt{test\_get\_current\_user\_decodes\_rs256} & Validação de assinatura RS256, \texttt{issuer} e \texttt{audience} \\ \hline
\end{tabularx}
\end{table}

\subsection{Validação Funcional}

O dashboard interativo permite validar visualmente todos os requisitos de segurança:
\begin{itemize}
  \item Aba Autenticação — login e decodificação de \gls{jwt}
  \item Aba Autorização — matriz de endpoints com resultados 200/403
  \item Aba JML — demonstração dos fluxos Joiner, Mover, Leaver
  \item Aba Auditoria — consulta de eventos com filtragem
  \item Aba MFA — fluxo de configuração e validação de \gls{totp}
\end{itemize}

\section{Considerações de Produção}

\subsection{Limitações da Validação Local}

A abordagem de validação local via cache \gls{jwks} tem uma limitação: se um token é revogado no Keycloak, a FastAPI não fica imediatamente ciente até que o token expire (300 segundos).

Esta limitação é compensada por:
\begin{itemize}
  \item Revogação explícita de sessões no Mover e Leaver
  \item Tempos de expiração curtos (5 minutos)
  \item Atualização periódica das chaves \gls{jwks}
\end{itemize}

\subsection{Escalabilidade}

A arquitetura é escalável:
\begin{itemize}
  \item Múltiplas instâncias da FastAPI podem validar tokens independentemente
  \item O Keycloak é o ponto centralizado de gestão de identidades
  \item A base de dados PostgreSQL é partilhada mas optimizada para auditoria
\end{itemize}

\subsection{Conformidade e Boas Práticas}

A solução segue as recomendações do NIST:
\begin{itemize}
  \item \textbf{NIST SP 800-63B} — AAL2 com \gls{mfa} para acesso a dados sensíveis
  \item \textbf{NIST SP 800-63C} — Federação de identidades via \gls{oidc}
  \item Princípio do privilégio mínimo — cada papel tem permissões específicas
  \item Princípio da defesa em profundidade — múltiplas camadas de validação
\end{itemize}
